\documentclass{article}
\usepackage[utf8]{inputenc}
\usepackage[polish]{babel}
\usepackage{graphicx}
\usepackage[T1]{fontenc}
\usepackage{wrapfig}
\usepackage{fancyhdr}
\usepackage{tabularx}




%\usepackage[a4paper,top=2cm,bottom=2cm,left=3cm,right=3cm,marginparwidth=1.75cm]{geometry}
\usepackage[includeheadfoot,
            left=1in,
            right=1in,
            top=1cm,
            bottom=1cm,
            headheight=1cm]{geometry}


\pagestyle{fancy}
\fancyhf{}  
\renewcommand{\headrulewidth}{0pt}  
\renewcommand{\footrulewidth}{0pt}  

\fancypagestyle{firstpage}
{
    \fancyhead[L]       % <-----------------LEWA STRONA-----------------
        {
     
            \vspace{0.25cm}            
            \textbf{\textit{Imię i nazwisko}} \\
            {\textit{Rok i kierunek studiów}} \\
            {\textit{Dzień, godzina}} \\
        }

        
    \fancyhead[R]   % <---------------------PRAWA STRONA------------
        {
         
             \vspace{0.25cm}
             \textit{Data}  \\ 
             {Prowadzący/a: } \\ 
             {\textit{tytuł imię nazwisko}} \\  
        
        }
}
\begin{document}
\textbf{ }\\
\textbf{ }\\
\thispagestyle{firstpage}
\centering
\section*{Ćwiczenie nr. }
\subsection*{Temat:}
\textbf{ }\\
\textbf{ }\\
\textbf{ }\\
\textbf{ }\\
\begin{tabularx}{0.8\textwidth} { 
  | >{\centering\arraybackslash}X |     % 1  }
  | >{\centering\arraybackslash}X |     % 2  }   LICZBA 
  | >{\centering\arraybackslash}X |     % 3  }   KOLUMN
  | >{\centering\arraybackslash}X |}    % 4  }
 \hline


 %---------------------------OPIS-------------------------------
 \# 
 & 1. wartość mierzona bezpośrednio [jednostka] 
 & 2. wartość mierzona bezpośrednio [jednostka] 
 & 3. wartość mierzona bezpośrednio [jednostka]  \\
%---------------------------OPIS-------------------------------

 
\hline
\hline
%---------------------------DANE-------------------------------
\hline 1 & 123 & 1 & 1\\
\hline 2 & 123 & 1 & 1 \\
\hline 3 & 123 & 1 & 1  \\
\hline 4 & 123 & 1 & 1  \\ 
\hline 5 & 123 & 1 & 1 \\
\hline 6 & 123 & 1 & 1 \\
\hline 7 & 123 & 1 & 1  \\
\hline 8 & 123 & 1 & 1  \\
\hline 9 & 123 & 1 & 1  \\ 
\hline 10 & 123 & 1 & 1 \\
\hline 11 & 123 & 1 & 1 \\
\hline 12 & 123  & 1 & 1 \\
\hline 13 & 123  & 1 & 1 \\
\hline 14 & 123  & 1 & 1 \\ 
\hline 15 & 123 & 1 & 1 \\
\hline 16 & 123 & 1 & 1 \\
\hline 17 & 123  & 1 & 1 \\
\hline 18 & 123  & 1 & 1 \\
\hline 19 & 123  & 1 & 1 \\ 
\hline 20 & 123 & 1 & 1 \\
\hline 21 & 123 & 1 & 1 \\
\hline 22 & 123  & 1 & 1 \\
\hline 23 & 123  & 1 & 1 \\
\hline 24 & 123  & 1 & 1 \\ 
\hline 25 & 123 & 1 & 1 \\
%---------------------------DANE-------------------------------
\hline
\end{tabularx}

\textbf{ }\\
\textbf{ }\\
\textbf{ }\\


\raggedright
    {
        {Dokładność wartości z pomiaru 1} \\
        {$\Delta$1= wielkość [jednostka]}\\
        \textbf{ }\\
        {Dokładność wartości z pomiaru 2} \\
        {$\Delta$2= wielkość [jednostka]}\\
        \textbf{ }\\
        {Dokładność wartości z pomiaru 3} \\
        {$\Delta$3= wielkość [jednostka]}\\
        \textbf{ }\\
    }

\pagebreak


\centering

\section*{ZAGADNIENIA TEORETYCZNE}

{We wstępie teoretycznym powinny być poruszone zagadnienia wymienione w instrukcji do
ćwiczenia oraz materiały źródłowe (BIBLIOGRAFIA), na podstawie których wstęp został
przygotowany.}

\section*{OPIS DOŚWIADCZENIA}

{W kilku zdaniach należy opisać poszczególne etapy wykonywania eksperymentu.}

\section*{OPRACOWANIE WYNIKÓW POMIARÓW}

{Należy wykonać wszystkie niezbędne obliczenia zawarte w instrukcji ćwiczenia.
Należy podać wzory, na podstawie których będą wykonywane obliczenia.
Jeżeli w obliczeniach wykorzystywane są wartości stałe należy podać ich wartości oraz
jednostki.
W przypadku wykonywania wielu obliczeń tego samego charakteru wystarczy przedstawić
kilka przykładowych wyliczeń i następnie w formie tabeli przedstawić wszystkie wyniki (wraz
z jednostkami).
W przypadku wykonywania obliczeń na komputerze należy podać nazwę wykorzystanego
programu. Jeżeli opracowanie wyników wymaga przygotowania wykresu należy wykonać go
w formacie A4.}\\



{Na podstawie instrukcji ONP oraz wskazówek zawartych w instrukcji danego ćwiczenia należy
obliczyć niepewności pomiarowe. Należy podać wzory, na podstawie których będą
wykonywane obliczenia. W przypadku wykonywania wielu obliczeń tego samego charakteru
wystarczy przedstawić kilka przykładowych wyliczeń i następnie w formie tabeli przedstawić
wszystkie wyniki.
Po obliczeniach, należy przedstawić wynik wraz z niepewnością z uwzględnieniem zasady 2
cyfr znaczących.}

\section*{WNIOSKI}

{Na podstawie wykonanych obliczeń i wyznaczonych wielkości fizycznych należy porównać
otrzymane wartości z wartościami tablicowymi. Należy również omówić oszacowane
niepewności pomiarowe (ich źródło, wpływ na wynik).}
\end{document}


%Dodatkowe uwagi:

1. Sprawozdanie może być pisane ręcznie. Proszę jednak o czytelność pisma!!!

2. Sprawozdanie MUSI zawierać wszystkie części (tabela pomiarową, teoria,
przebieg ćwiczenia, obliczenia, niepewności, wnioski i wykresy). Brak
jakiejkolwiek części kwalifikuje do zwrotu złożonego sprawozdania bez dalszego
sprawdzania.

3. Wykresy należy zamieszczać na osobnych kartkach (format A4). Wykonywać za
pomocą komputera lub ręcznie na papierze milimetrowym. Należy tak dobrać
skalę, aby wykres zajmował całą stronę.

4. Punktów pomiarowych naniesionych na wykresach nie łączymy! W przypadku
dopasowania prostej regresji, wraz punktami na wykresie należy nanieść prostą
regresji.

5. Na wykresach razem z punktami należy nanieść niepewności pomiarowe w formie
tzw. krzyży niepewności pomiarowych.

6. Do sprawozdania należy dołączyć kartkę pomiarową z ćwiczenia podpisaną przez
prowadzącego.

7. Przy zapisie wyników wraz z niepewnością obowiązuje zasada podawania 2 cyfr
znaczących (instrukcja ONP).

8. Niepewności pomiarowe w większości przypadków wyliczamy bazując na trzech
metodach:
a) gdy mamy pomiary skorelowane korzystamy z zależności 17 w instrukcji ONP,
b) gdy mamy pomiary nieskorelowane korzystamy z zależności 15 w ONP,
c) w przypadku dopasowywania prostych regresji, niepewności obliczamy ze
wzorów 6 i 7 w ONP.